\section{Next Work}

\paragraph{Week 3 selection.}
The packed prefix closed, while the parent raw call-site \(\mu\)-hash claim is
still partial. The decision rule therefore selects another focused
SHAKE/control/memory sprint:
\begin{enumerate}
  \item prove an exact refinement between
  \procname{\_sf\_mu\_rawpre}, the raw branch of
  \procname{\_\_verify\_hash\_mu}, and the compiled raw wrappers;
  \item lift the constructed 992-byte relation through actual Sign/Verify
  callers and public memory contracts; and
  \item rerun extraction hashes, fresh EasyCrypt compilation, non-vacuity
  checks, baselines, and the LaTeX build.
\end{enumerate}

\paragraph{Go/no-go.}
This is a \emph{go} for the narrow control-flow and memory-lifting work. It is
a \emph{no-go} for widening the claimed public-API refinement or paper-level
EUF-CMA result. Highbits/LSB packing and full challenge-call equivalence begin
only after \claimid{OBL-MU-TOPLEVEL-CONTROL} closes.

\paragraph{Minimum next theorem.}
The next nontrivial deliverable must show that an actual terminating raw
generated Sign/Verify call establishes the already-proved wrapper pre/post
relation. It must expose branch and memory premises and must not introduce a
SHAKE or implementation-semantic axiom.

\paragraph{Longer path.}
After the exact raw call site closes, the order is public-context encoding,
highbits/LSB equality, full challenge/ROM query equivalence, then the broader
Sign/Verify arithmetic and security-distribution obligations. KeyGen
termination, FFT/TieMin, sampler/rejection distributions, and forking remain
separate workstreams.
